\section{全球化、新加坡与“跑路”误解}

访谈最后，Peak 谈到新加坡。他的回答很务实：Manus 选择新加坡，不是因为“跑路”，而是因为公司从一开始服务全球市场，需要一个能支持国际客户、合规、团队集中办公和全球业务的地点。这个问题和 Agent 产品有关：如果用户、客户和支付都在全球，团队就必须提前处理合规、数据、支付、法律和客户信任。

\subsection{为什么不是简单地理选择}

前面说新加坡不是情绪化选择，本节拆开实际原因。Peak 提到 Manus 已经通过 SOC 2、ISO 27001 / 27701、GDPR 等合规工作。这类工作对普通 C 端工具可能显得很重，但对 Agent 公司很关键。因为 Agent 可能接触用户账户、文件、浏览器、企业资料、支付和自动化操作。全球客户愿不愿意使用，往往取决于安全合规能否先被信任。

\begin{knowledgebox}{全球化 Agent 公司的合规压力}
\begin{tabularx}{\textwidth}{p{0.22\textwidth}p{0.34\textwidth}X}
\toprule
合规/治理点 & 要回答的问题 & 和 Agent 的关系 \\
\midrule
数据保护 & 用户数据在哪里、如何处理 & Agent 需要大量上下文和文件。 \\
安全审计 & 系统是否有控制和追踪 & 自动化操作必须可追责。 \\
跨境服务 & 客户在哪个市场，如何签约 & 全球化产品要适配不同规则。 \\
团队协作 & 团队是否集中，沟通成本如何 & 早期团队要快速迭代和决策。 \\
\bottomrule
\end{tabularx}
\end{knowledgebox}

\subsection{“跑路”为什么是错误表述}

所谓“跑路”把全球化、合规和客户市场误解成地理逃离。Peak 的解释是，Manus 一直做全球市场，新加坡是为了服务客户和组织团队。对 AI 公司来说，客户在哪里、合规要求在哪里、支付和法律结构在哪里，办公室和总部选择就会跟着变化。

\begin{warningbox}{不要把全球化简化成情绪词}
AI 公司选择海外节点，可能来自客户、合规、融资、支付、人才和团队协作等多重原因。用“跑路”概括，会遮蔽真正的商业和制度问题。
\end{warningbox}

\subsection{本章小结}

新加坡选择说明，Agent 公司不是纯线上产品团队。它需要面对全球客户、合规制度、数据安全和组织协作。这些问题同样属于“AI 像制造业”的现实部分。

